\chapter{Approfondir les preuves et préparer la suite}
\section*{Exercice 24.1 -- L'ordre des quantificateurs}
\textbf{Facultatif.} Comparer les affirmations : « pour tout réel $x$, il existe un réel $y$ tel que $y>x$ » et « il existe un réel $y$ tel que, pour tout réel $x$, $y>x$ ». Écrire leurs négations en français.

\section*{Exercice 24.2 -- Une somme télescopique}
\textbf{Facultatif.} Démontrer $1+3+\cdots+(2n-1)=n^2$ en utilisant $k^2-(k-1)^2$. Vérifier pour $n=4$ puis expliquer pourquoi cette vérification ne suffit pas.

\section*{Exercice 24.3 -- Une récurrence guidée}
\textbf{Facultatif.} Pour $u\ge0$ et $n\in\N$, démontrer par récurrence $(1+u)^n\ge1+nu$. En déduire que pour tout seuil $M>0$ et $u>0$, les puissances $(1+u)^n$ finissent par dépasser $M$.

\section*{Exercice 24.4 -- Une borne de fluctuation}
\textbf{Facultatif.} Pour une fréquence $F_n$ de succès lors d'épreuves indépendantes de probabilité $p$, le cours démontre $P(|F_n-p|\ge\varepsilon)\le1/(4n\varepsilon^2)$. Calculer cette borne pour $n=10\,000$ et $\varepsilon=0{,}05$. Peut-on en déduire que tout échantillon est proche de $p$ ?

\section*{Exercice 24.5 -- Trouver le trou d'une preuve}
\textbf{Facultatif.} Un texte dit : « $a=b$, donc $a^2=ab$, donc $(a-b)(a+b)=b(a-b)$, donc $a+b=b$, donc $2=1$ en prenant $a=b=1$ ». Localiser et expliquer l'erreur.

\section*{Exercice 24.6 -- Démontrer seul}
\textbf{Facultatif.} Soient $a,b>0$. Démontrer $\frac{a+b}{2}\ge\sqrt{ab}$, préciser le cas d'égalité, puis en déduire $\frac1a+\frac1b\ge\frac4{a+b}$.
